\chapter{Daily use}

\label{ch:daily}

A normal day with pii is short. You ask the daemon what it is, you give it work --- a
prompt, a graph, a queue row --- and then you read what it did. This chapter is that loop
in the order you actually run it, with the exact flags the binary accepts at the revision
this manual describes, the exit code each command returns, and the places where the answer
looks like a success and is not. Nothing here needs the repository; everything here assumes
a daemon that is already running, as installed in Chapter~\ref{ch:install}.

\begin{center}
\resizebox{\linewidth}{!}{
\begin{tikzpicture}[node distance=9mm and 12mm]
  \node[b] (health) {\texttt{pii health}\\\scriptsize is it the build I think it is?};
  \node[b, right=of health] (work) {give it work\\\scriptsize prompt, graph, queue row};
  \node[b, right=of work] (read) {read the run\\\scriptsize get, trace, explain};
  \node[b, above=of read] (keep) {keep the stores\\\scriptsize stats, compact, prune};
  \node[b, left=of keep] (fine) {as expected?};
  \node[b, left=of fine] (report) {report it\\\scriptsize the journal is the evidence};
  \draw[f] (health) -- (work);
  \draw[f] (work) -- (read);
  \draw[f] (read) -- (keep);
  \draw[f] (keep) -- (fine);
  \draw[bk] (fine) -- node[p,above=2pt,midway]{no} (report);
  \draw[f] (fine) -- node[p,right=2pt,midway]{yes} (health);
\end{tikzpicture}}
\captionof{figure}{The daily loop. The first and last steps are the ones that get skipped: a
health call before you trust a result, and a store reading before the disk tells you.}
\end{center}

\section{Asking the daemon a question}

\texttt{pii health} is the first command of any session, including a session that starts
with \enquote{it worked yesterday}. It prints the daemon's own account of itself and
returns an exit code that a script can branch on.

\begin{lstlisting}[language=bash]
pii health                 # human text; 0 = same build, 2 = no daemon, 3 = drift
pii health --json          # the sys.status payload, verbatim
pii metrics                # one name-value line per metric
pii disk                   # artefact roots with size, age and filesystem
\end{lstlisting}

The text form is documented line by line in Chapter~\ref{ch:what} and is not repeated
here. The useful habit is narrower: read \texttt{client:} against the first line, because
that pair is the only thing in pii that tells you whether the binary in your shell and the
process answering you are the same program. A \texttt{DRIFT} line on stderr with exit 3 is
not an error to fix; it is a fact to act on --- restart the service, or use the binary the
daemon is running.

\texttt{pii disk} answers a different question, and it exists because the answer used to
take four commands and a guess. It walks a fixed list of artefact roots
(\texttt{\_tmp/gate-target}, \texttt{\_tmp/gate-target-fresh}, \texttt{\_tmp/targets},
\texttt{\_tmp/\_trash}, \texttt{/projects/\_tmp}, \texttt{/projects/\_trash},
\texttt{/tmp/\_trash-podman}, \texttt{/tmp/gate-cache} and every unclaimed
\texttt{\_tmp/*target*} directory), prints each one's size, age, filesystem and
whether an owner claims it, then prints free space for \texttt{/} and for the volume that
holds the repository. It is read-only; removal is \texttt{cargo xtask disk --gc}, which
owns the TTL.

The reason it prints \emph{two} free-space lines is a measured defect: the gate's floor
was evaluated on \texttt{/} while the artefacts landed on \texttt{/projects}, so a host
with 16\,GB free on a 288\,GB overlay and 147\,GB free on the repository volume refused to
run --- and the obvious first move, \texttt{cargo clean} on the repository, freed space on
the volume that was never short.

\section{One-shot prompts}

\texttt{pii -p} is the shortest path to a model. It carries a prompt, and its flags decide
which model, which graph, which workspace and how much confinement.

\begin{lstlisting}[language=bash]
pii -p "summarise the open items in specs/350-an-incident-becomes-a-case.md"
pii -p "why did run r000000000130 end unverified" --model llamacpp-server:granite-4.2-3b
pii -p "check the tree" --graph preflight --workspace /projects/pi-improved
\end{lstlisting}

\begin{longtable}{@{}p{0.20\textwidth}p{0.72\textwidth}@{}}
	\toprule
	\textbf{Flag} & \textbf{Effect} \\
	\midrule
	\endhead
	\texttt{--model M} & the model spec for this call; wins over \texttt{\$PII\_MODEL} and
	over \texttt{[core].model} \\
	\texttt{--graph G} & run an installed graph instead of the default one-shot path \\
	\texttt{--workspace D} & the directory the sandbox and the tools are scoped to;
	defaults to the current directory \\
	\texttt{--sandbox off|workspace} & \texttt{workspace} is the default; \texttt{off}
	runs unconfined and says so \\
        \texttt{--connect auto|loopback|\textless{}socket\textgreater{}} & which transport to
use; \texttt{auto} takes
	a live daemon socket if one answers \\
	\texttt{--json} & machine-readable output; already the default when stdout is not a
	terminal \\
	\texttt{--input J} & extra graph input, merged into the run input as JSON \\
	\bottomrule
\end{longtable}

Two behaviours are easy to observe and worth knowing before the first surprise. The
transport is chosen for you: a live daemon socket wins, and only if nothing answers does
the command build the engine in its own process. The daemon is never started for you ---
there is no auto-start anywhere in the path --- so a \texttt{-p} on a host whose service is
down silently becomes an in-process engine with its own view of the state directory. And
if no model resolves at all, the command exits 5 with one line:
\enquote{no model configured (set PII\_MODEL or --model)}. That is a different failure from
a model the provider does not serve, which fails later, at spawn.

\section{Running a graph}

There are two ways to run a graph, and they are not interchangeable. By \emph{name}, you
ask the daemon, which serves the installed set and owns the stores. By \emph{file}, you run
it yourself against a store you name, which is what the test suites and a one-off
diagnostic do.

\begin{lstlisting}[language=bash]
# by name: the daemon resolves it, runs it, and journals it
pii run start queue --input '{"workdir":"/projects/pi-improved","base":"<40-hex>"}'

# by file: you own the store path and the write grants
pii graph run graphs/preflight.toml --db ~/.local/state/pii/store.redb \
  --input '{}' --allow-write /projects/pi-improved/_tmp
\end{lstlisting}

\begin{center}
\resizebox{\linewidth}{!}{
\begin{tikzpicture}[node distance=9mm and 12mm]
  \node[b] (name) {\texttt{pii run start \textless{}graph\textgreater{}}};
  \node[b, right=of name] (daemon) {daemon\\\scriptsize installed graphs, its stores};
  \node[b, right=of daemon] (journal) {journal\\\scriptsize run id, every port};
  \node[b, below=of name] (file) {\texttt{pii graph run \textless{}file-> --db}};
  \node[b, right=of file] (local) {this process\\\scriptsize the engine, in-process};
  \node[b, right=of local] (journal2) {journal in \texttt{--db}\\\scriptsize same format};
  \draw[f] (name) -- (daemon);
  \draw[f] (file) -- (local);
  \draw[f] (daemon) -- (journal);
  \draw[f] (local) -- (journal2);
  \draw[sf] (name) -- node[p,left=2pt,midway]{no daemon?} (file);
\end{tikzpicture}}
\captionof{figure}{Two paths, one journal format. The name path is what an installation uses;
the file path is what a diagnostic uses, and it is the only one where you choose the store.}
\end{center}

The offline path builds its own capability root, and the defaults are worth memorising
because they explain most \enquote{permission denied} results from a hand-run graph:
reads are allowed for the current directory, writes are allowed for
\texttt{\textless{}cwd\textgreater{}/\_tmp} plus every \texttt{--allow-write} directory, and the
only process
argv[0] permitted is \texttt{/bin/sh}. Everything else is denied by the capability check
before the node runs. \texttt{--allow-write} is canonicalised, so a relative or symlinked
grant would otherwise grant nothing.

Both paths share one exit-code vocabulary, and the two codes that are not 0 or 1 are the
interesting ones:

\begin{center}
\begin{tabular}{@{}p{0.10\textwidth}p{0.78\textwidth}@{}}
	\toprule
	\textbf{Code} & \textbf{Meaning} \\
	\midrule
	0 & the run reached its terminal node and reported success \\
	1 & the run failed, or could not be started; the message names the reason \\
	2 & usage error: an unknown flag, a missing value, \texttt{--input} that is not JSON \\
	3 & the validator refused the graph --- and, on this path, no store file was created \\
	4 & the run suspended and holds a resume key \\
	\bottomrule
\end{tabular}
\end{center}

The distinction between 3 and 1 is the reason a refused graph creates no store file: the
graph is loaded and validated \emph{before} the store is opened, so a wrong pin or a
missing callee costs nothing and leaves nothing behind.

\section{The queue}

The queue is how pii is given work rather than a question. A row carries an id, a title, a
spec path, the criteria that decide whether the work is done, and the paths the row expects
to touch. Every subcommand is one protocol call; none opens the store.

\begin{lstlisting}[language=bash]
pii queue add --id 376 --title "retention drops what the policy expires" \
  --spec /projects/pi-improved/specs/376-retention-drops-what-the-policy-expires.md \
  --criteria "cargo test -q -p pii-store --test retention_drops_what_the_policy_expires" \
  --touches crates/pii-store/src/retention.rs
pii queue ls
pii queue claim --by operator
\end{lstlisting}

The spec path is resolved by the daemon, not by your shell. A relative path is the single
most common first mistake with \texttt{queue add}: the daemon's working directory is the
unit's \texttt{WorkingDirectory}, which is the state directory, so a relative
\texttt{specs/376-*.md} resolves to a file that does not exist. Use an absolute path and
the row will be claimable.

Running the queue is one graph run, and the graph needs two inputs: the work directory and
the commit that holds the row's spec. Both are passed as one JSON object to
\texttt{pii run start queue --input}, with \texttt{workdir} an absolute directory and
\texttt{base} a forty-hex commit. The spec graph's own limits are shared by the call, so
\texttt{max\_wall\_ms} on the queue's \texttt{[limits]} table is the whole run's wall, not
each step's.

\begin{center}
\resizebox{\linewidth}{!}{
\begin{tikzpicture}[node distance=8mm and 12mm]
  \node[b] (select) {\texttt{select}\\\scriptsize claims one row};
  \node[b, right=of select] (impl) {\texttt{implement}\\\scriptsize calls the spec graph};
  \node[b, right=of impl] (record) {\texttt{record}\\\scriptsize writes the verdict};
  \node[b, below=of select] (none) {\texttt{record}\\\scriptsize empty or blocked};
  \node[s, below=of impl] (store) {queue store\\\scriptsize the row's status};
  \draw[f] (select) -- node[p,above=2pt,midway]{ok} (impl);
  \draw[f] (impl) -- node[p,above=2pt,midway]{ok} (record);
  \draw[f] (select) -- node[p,left=2pt,midway]{empty} (none);
  \draw[bk] (impl.north) to[bend left=22] node[p,above=2pt,midway]{refused} (record.north);
  \draw[sf] (record.south) -- (store.north);
\end{tikzpicture}}
\captionof{figure}{\texttt{queue} claims one row, hands it to the implementer, and records a
verdict. The red edge is the one that took a row from \enquote{stuck Running} to
\enquote{recorded as refused}.}
\end{center}

The three ports of \texttt{select} are the whole reason the queue is readable.
\texttt{ok} means a row was claimed, \texttt{empty} means there was nothing to do, and
\texttt{blocked} means a row exists but something it waits on does not --- and \texttt{blocked}
ends the run on its own edge rather than pretending to have worked. A row that stays
\texttt{Running} after the run ends is the failure this shape exists to prevent: at the
revision this manual describes, a refused callee used to dead-end inside the splice
because the callee's refusal nodes declare no \texttt{ok} port, so no edge fired and the
row was never recorded. The fix wires the refusal by name ---
\texttt{implement.refused}, \texttt{implement.unverified} --- and the queue's own graph now
carries both edges.

The rule for an operator is one sentence: if a row does not leave \texttt{Running}, read
the trace and look at the last node, not at the last message. The verdict is in the
journal even when the record was never written.

\section{Reading a run}

Three commands read a run, and they answer different questions.

\begin{lstlisting}[language=bash]
pii run ls --status suspended         # running|suspended|succeeded|failed
pii run ls --graph queue              # filter by graph name, client-side
pii run get <run_id>                  # the header: graph, status, timing
pii run trace <run_id>                # node by node, with the port that fired
pii run explain <run_id>              # the same journal, summarised
\end{lstlisting}

\texttt{--status} accepts exactly four names and refuses anything else with the list in
the message. \texttt{refused} is a run status the store knows but not one the filter
accepts, which is a small asymmetry worth knowing before you script it.

\texttt{pii run get} distinguishes three answers that look alike in a hurry. A run the
daemon holds but has not journalled yet answers \texttt{accepted} --- the id exists and the
journal has not materialised, so you poll again. A stored state answers
\texttt{running}, \texttt{suspended}, \texttt{succeeded} or \texttt{failed}. An id the
daemon has never heard of answers \enquote{unknown to this daemon}: it was never issued, or
the daemon restarted and lost an accepted-but-never-journalled start. Polling an
\texttt{accepted} id until it becomes terminal is correct; treating the first
\texttt{not found} as a failure is not, and the measured baseline is that a validated start
never polls \texttt{not found} for ever --- a start the run's own validation refuses still
returns an id, then goes to \texttt{failed} with the reason in the trace.

\begin{center}
\resizebox{\linewidth}{!}{
\begin{tikzpicture}[node distance=13mm and 16mm]
  \node[b] (start) {\texttt{run start}};
  \node[b, right=of start] (acc) {\texttt{accepted}\\\scriptsize no journal yet};
  \node[v, right=of acc] (state) {stored?};
  \node[b, right=of state, yshift=14mm] (run) {\texttt{running}};
  \node[b, right=of state, yshift=-14mm] (done) {\texttt{succeeded} / \texttt{failed}};
  \node[b, below=of state] (unknown) {unknown\\\scriptsize never issued, or lost};
  \draw[f] (start) -- (acc);
  \draw[f] (acc) -- node[p,above=2pt,midway]{poll} (state);
  \draw[f] (state) -- node[p,above=2pt,midway]{live} (run);
  \draw[f] (state) -- node[p,below=2pt,midway]{terminal} (done);
  \draw[bk] (state) -- node[p,right=2pt,midway]{no row} (unknown);
\end{tikzpicture}}
\captionof{figure}{Three answers from \texttt{run get}, and only one of them is a polling state.
The red edge is the reading that is a real problem: an id the daemon never issued, or lost
across a restart before its journal was written.}
\end{center}

Read a trace from the end. The last \texttt{NodeFinished} names the node and the port it
took, and that pair is the decision the run made. A run that stopped without reaching a
terminal node has not failed loudly --- it has failed quietly, which is the failure mode
the journal exists to make visible. The events you will read most are
\texttt{Started}, \texttt{TickStarted}, \texttt{NodeStarted}, \texttt{NodeFinished},
\texttt{NodeFailed}, \texttt{NodeProgress}, \texttt{NodeCancelled}, \texttt{NodeDead},
\texttt{NodeInterrupted}, \texttt{Suspended}, \texttt{Resumed} and \texttt{NodeObserved};
Chapter~\ref{ch:what} lists all seventeen.

\texttt{NodeProgress} deserves its own sentence, because it is the only event that
separates a slow node from a wedged one: it carries elapsed milliseconds and the child's CPU
milliseconds, so \emph{elapsed grows, CPU flat} is a node waiting on something that will
never answer, while \emph{both grow} is a node doing work.

\subsection*{The offline trace, and its queries}

When the daemon is down, the same journal is readable from the store file directly, and
the offline form answers more questions than the online one.

\begin{lstlisting}[language=bash]
pii graph trace <run_id> --db ~/.local/state/pii/store.redb
pii graph trace <run_id> --db ... --terminal
pii graph trace <run_id> --db ... --node-status verify
pii graph trace <run_id> --db ... --reason refused
pii graph trace <run_id> --db ... --overlapped code,docs
\end{lstlisting}

With no query at all, the command prints the journal as NDJSON: one event per line, each
with a sequence number.
\texttt{--terminal} prints the run's terminal value, \texttt{--node-status} prints one
node's last port, \texttt{--reason} prints why a refusal node refused,
\texttt{--max-batch-width} reports how wide the scheduler's batches ran, and
\texttt{--overlapped a,b} answers whether two nodes were ever in flight at the same time ---
which is the question you ask when a graph is suspected of hiding a serialisation.

\subsection*{When a run needs a hand}

Four commands exist for runs that will not finish on their own, and they differ in what
they do to the evidence.

\begin{center}
\begin{tabular}{@{}p{0.26\textwidth}p{0.62\textwidth}@{}}
	\toprule
	\textbf{Command} & \textbf{What it does} \\
	\midrule
        \texttt{pii run resume \textless{}id\textgreater{}} & continues a \texttt{Suspended} run
through its resume
	key \\
        \texttt{pii run retry \textless{}id-> --node \textless{}id\textgreater{}} &
re-runs one node of a finished run \\
        \texttt{pii run abandon \textless{}id-> --reason T} & marks the run
abandoned, with the reason
	recorded \\
        \texttt{pii run cancel \textless{}id-> [--kill]} & cancels a live run;
\texttt{--kill} also kills
	the child \\
	\bottomrule
\end{tabular}
\end{center}

Abandon and cancel are the two that change what retention will later do. A run left
\texttt{Running} or \texttt{Suspended} is never swept, however old it is, so a run you
have given up on stays in the store for ever unless you end it deliberately.

\section{Stores, memory and the indices}

Store maintenance is three commands, and two of them need the daemon stopped.
\texttt{pii store stats} prints one line per table with its row count, then the file's
byte size. \texttt{pii store compact} takes redb's exclusive lock, so a running daemon makes
it refuse rather than copy a file mid-write, and it prints the measured size of each file
before and after. \texttt{pii store prune} takes the same lock, runs one retention sweep and
prints \texttt{pruned N rows}.

\texttt{pii store stats} is the command to run weekly. It prints
\texttt{store.redb \textless{}table-> rows=N} for every table and then
\texttt{store.redb bytes=N}, which together are the only honest growth signal --- the file
size alone is not, because redb never shrinks and a store that has been compacted tells you
nothing about how fast it is filling. \texttt{pii store compact} prints
\texttt{compacted \textless{}name-> \textless{}before-> ->
\textless{}after-> bytes} per file, and a store with nothing to
reclaim prints a near-equal pair, which is the truth rather than a failure.

\texttt{pii store prune} runs one retention sweep under the default policy and prints the
count. The defaults are thirty days for finished runs, seven for node-call events and at
most one thousand rows per sweep; the sweep never removes a \texttt{Running} or
\texttt{Suspended} run, and it keeps any row whose creation time is missing. The daemon
also runs one sweep at boot, so a host that restarts nightly trims itself.

Memory has its own command group, and the two that are used daily are search and the
entity graph: \texttt{pii memory search \textless{}q\textgreater{}} with an optional
\texttt{--k} (six results by
default), and \texttt{pii memory put \textless{}text\textgreater{}} with \texttt{--layer} and
\texttt{--kind} for
a document's tier. \texttt{pii memory tiers} says what sits in each tier,
\texttt{pii memory entity list} gives the entities and their counts, and
\texttt{pii memory neighbours} gives one hop around an entity. The rest of the group
(\texttt{learn}, \texttt{consolidate}, \texttt{link}, \texttt{compact}, \texttt{archive},
\texttt{restore}) is maintenance, and each one either refuses or reports when its
precondition is missing rather than doing half the work.

\texttt{pii index code} and \texttt{pii index docs} fill the code and documentation indices
from a tree, and \texttt{pii skill list} answers what the daemon knows how to do. Both are
read-mostly and safe to re-run: an index command re-reads and updates, it does not
duplicate.

\section{A cadence that works}

The commands above are cheap; the discipline is knowing which ones to run on a clock. The
cadence below is the one Chapter~\ref{ch:operate} assumes, written as a table so it can be
checked rather than remembered.

\begin{longtable}{@{}p{0.13\textwidth}p{0.30\textwidth}p{0.45\textwidth}@{}}
	\toprule
	\textbf{When} & \textbf{Command} & \textbf{What you are looking for} \\
	\midrule
	\endhead
	Every session & \texttt{pii health} & exit 0; \texttt{client:} equals the build line \\
	Every session & \texttt{pii run ls --status running} & a run that should have finished \\
	Daily & the log file, last lines & a log that ends on \texttt{daemon.started} means a
	hard death \\
	Weekly & \texttt{pii store stats} & rows and bytes per table, recorded, so the slope
	exists \\
	Weekly & \texttt{pii disk} & an artefact root growing without an owner \\
	Monthly & \texttt{pii store compact} & the tail slack back; nothing more \\
	Monthly & a restore into a scratch directory & a backup nobody restored is a hope \\
	\bottomrule
\end{longtable}

Two of those rows are worth the cost even on a quiet host. A weekly \texttt{store stats}
sample is what turns \enquote{the disk is filling} into a rate, and a monthly restore into a
scratch state directory is what turns a backup into a backup --- both drills are described
in Chapter~\ref{ch:operate}.

\section{What a quiet day looks like}

A host with nothing to do is not silent, and knowing the small noises saves a false alarm.
The daemon writes a \texttt{daemon.started} event, keeps its health record current, and
runs one retention sweep at boot. If a channel is configured it polls on its interval; if a
worker is configured it says hello once. Nothing else should appear in the log.

The four readings that answer \enquote{is this host doing anything} are
\texttt{pii health} (first three lines: build, client, config hash), \texttt{pii run ls}
(usually an empty list on a host with no work), \texttt{pii cron ls} (the schedules the
daemon knows) and the last three lines of \texttt{\$STATE\_DIR/pii.log}. On a healthy host
the log's last line is a \texttt{daemon.started} only after a restart --- a log that ends on
\texttt{daemon.started} while the process has been up for days means the daemon died hard
and was restarted by the unit, and \texttt{since\_deploy.hard\_deaths} in \texttt{pii
health} will say how many times.

If \texttt{pii run ls} is empty and you expected rows, the question is not \enquote{where did
the runs go} but \enquote{which store is this daemon writing} --- and \texttt{pii health}
answers it with a path. That single habit, reading the store path before believing an empty
list, resolves more \enquote{lost work} reports than any other diagnostic in this manual.
